\section{So How Far Are We? An Honest Distance Estimate}
\label{sec:distance}

We can now answer the title question without evasion. The answer is not a
number of years but a map: the distance depends entirely on which rung of
the ladder one means, and the rungs differ not by degree but by kind.
Table~\ref{tab:distance} gives our estimate, with the evidence and the
dominant blocker for each.

\subsection{Rung 1--2: essentially here, and useful}
Signal detection and affect classification are solved in the laboratory
and shipping in products. The remaining work is engineering and
validation---public benchmarks, calibration, multimodal fusion---not
discovery. A calibrated, validated welfare monitor that reports arousal and
valence with honest uncertainty is plausibly a \textbf{2--4 year} horizon
from the research frontier to reliable deployment. This is the genuinely
useful payoff: early stress or pain detection, shelter welfare assessment,
veterinary triage. It is not ``translation,'' and calling it such would
oversell it; but it is real.

\subsection{Rung 3: a foothold, a decade of work}
Intent and context inference on \emph{natural} signals is emerging. The
soundboard literature proves that referential intent exists and is
checkable~\cite{bastos2024plos,bastos2024scirep}, but it does so with
human-supplied symbols. Extending this to spontaneous barks, faces, and
postures---fusing modalities, modeling goals, validating closed-loop---is a
research program, not a product cycle. We estimate \textbf{validated
natural-signal intent inference at roughly 2030 or later}, contingent on
the evaluation infrastructure of Section~\ref{sec:eval} being built. The
blocker is not model capacity; it is labeled, grounded, validated data.

\subsection{Rung 4--5: possibly never, and that is informative}
Open referential translation and free dialogue presuppose a generative
symbolic system in the dog that the biological evidence does not support.
Here the honest answer is \textbf{no credible timeline}, and the reason is
not pessimism about AI but a fact about the subject: there may be no
language there to translate. This is the survey's central, deflationary
claim. It does not rule out an \emph{engineered} dialogue---structured
human-to-dog input via trained symbols, with the dog responding through a
soundboard---which is a Rung-3 artifact wearing Rung-5 clothing. But faithful
translation of spontaneous canine ``speech'' into English sentences is, in
the strong sense, possibly ill-posed, and more compute will not fix a
category error.

\subsection{Why ``more scale'' is not the answer}
It is tempting to read the cross-species successes as evidence that scaling
foundation models will eventually crack the dog. The reading is mistaken.
CETI, DolphinGemma, and the finch models advance \emph{structure discovery
and generation}---the front of the pipeline. They do not supply semantics,
and their authors say so~\cite{sharma2024whale,google2025dolphingemma}. The
canine bottleneck is at the \emph{back} of the pipeline (grounding,
validation) and, above Rung 3, in the subject itself. Scaling the encoder
makes Rung 1--2 better; it does not create a grammar that is not there.
The distance to Rung 5 is therefore not a function of parameters but of
biology.

\subsection{The reframe: from translation to interpretation}
The productive move is to stop asking ``when will we translate dogs?'' and
start asking ``how well can we \emph{interpret} canine state, with calibrated
uncertainty and animal-grounded validation?'' That question has answers,
has metrics, and has welfare value. It also respects the animal: it treats
the dog as a signaler with its own modalities, not as a defective human
whose sentences we are waiting to recover. The Rosetta Stone that titles
this survey is, in the end, a warning as much as a promise: there may be no
decree to transcribe (Section~\ref{sec:conclusion}).

\begin{table}[htbp]
\centering
\caption{Distance estimate by rung. Horizons are judgments from the
evidence reviewed, not predictions; ``blocker'' names the dominant
obstacle. Rungs 4--5 are marked possibly ill-posed because the requisite
generative symbolic system is not evidenced in natural canine communication.}
\label{tab:distance}
\small
\begin{tabular}{@{}p{2.7cm}p{2.5cm}p{4.0cm}p{3.6cm}@{}}
\toprule
\textbf{Rung} & \textbf{Status} & \textbf{Horizon} & \textbf{Dominant blocker} \\
\midrule
1 Signal & Solved & Now & Engineering only \\
2 Affect & Shipping & 2--4 yr to validated welfare use &
Calibration, public benchmarks \\
3 Intent & Foothold & $\sim$2030+ & Grounded, validated data \\
4 Reference & Unproven & No credible timeline & No productive symbol system \\
5 Dialogue & Possibly ill-posed & Possibly never & No evidenced generative grammar \\
\bottomrule
\end{tabular}
\end{table}